\bookchapter{Κεφάλαιο 45: Διασφάλιση ποιότητας — κάνοντας τον κώδικα της AI αξιόπιστο}

\chapterintro{Ο κώδικας που γράφει η AI έχει ένα χαρακτηριστικό: συνήθως «φαίνεται πολύ σωστός» — τυπική ονο$\mu$ατολογία $\mu$εταβλητών, πλήρη σχόλια, καθαρή λογική. Αλλά το «φαίνεται σωστός» δεν είναι το ίδιο $\mu$ε το «είναι πραγ$\mu$ατικά σωστός». Τα κρυ$\mu$$\mu$ένα προβλή$\mu$ατα στον κώδικα AI συχνά είναι πιο δύσκολο να εντοπιστούν απ' ό, τι σε κώδικα γρα$\mu$$\mu$ένο $\mu$ε το χέρι, επειδή η επιφανειακή του ποιότητα είναι τόσο υψηλή που σε κάνει να χαλαρώσεις την προσοχή σου. Αυτό το κεφάλαιο σου δείχνει πώς να χτίσεις ένα αποτελεσ$\mu$ατικό σύστη$\mu$α διασφάλισης ποιότητας για τον κώδικα AI — $\mu$ε τρεις γρα$\mu$$\mu$ές ά$\mu$υνας, τα tests, την επιθεώρηση και τον έλεγχο εκδόσεων, ώστε το «φαίνεται σωστός» να γίνει «είναι πραγ$\mu$ατικά σωστός».}

\booksection{Ανάπτυξη $\mu$ε οδηγό τα tests: άσε την AI να γράψει τα tests}

Τα tests είναι η πρώτη γρα$\mu$$\mu$ή ά$\mu$υνας της ποιότητας στον προγρα$\mu$$\mu$ατισ$\mu$ό $\mu$ε AI. Τα καλά νέα: το να ζητήσεις από την AI να γράψει tests είναι πιο εύκολο και πιο αξιόπιστο από το να γράψει κώδικα λειτουργιών — επειδή η ανα$\mu$ενό$\mu$ενη συ$\mu$περιφορά ενός test είναι σαφής, δεν παράγεται εύκολα κώδικας «που φαίνεται σωστός αλλά στην πραγ$\mu$ατικότητα δεν είναι».

Η ροή του προγρα$\mu$$\mu$ατισ$\mu$ού $\mu$ε AI $\mu$ε οδηγό τα tests: πρώτα ζητάς από την AI τον κώδικα των λειτουργιών· όταν επαληθεύσεις ότι οι βασικές λειτουργίες δουλεύουν, της ζητάς α$\mu$έσως test cases. Τα test cases πρέπει να καλύπτουν την κανονική πορεία, τις οριακές περιπτώσεις και τις εξαιρέσεις. Μετά ζητάς από την AI να διορθώσει τον κώδικα των λειτουργιών $\mu$ε βάση την ανατροφοδότηση των tests. Το πλεονέκτη$\mu$α αυτής της ροής: $\mu$όλις γραφτούν, τα tests γίνονται διαρκές εργαλείο επαλήθευσης — κάθε $\mu$ελλοντική αλλαγή $\mu$πορεί να ελεγχθεί τρέχοντας τα tests, ώστε να επιβεβαιώνεις ότι δεν $\mu$πήκαν προβλή$\mu$ατα παλινδρό$\mu$ησης (regression).

\begin{lstlisting}
Γράψε unit tests για τη συνάρτηση calculate_discount.
Τα τεστ πρέπει να καλύπτουν:
1. Κανονική ροή:
   - προϊόν 100 γιουάν $\mu$ε 20% έκπτωση -> 80
   - προϊόν 200 γιουάν $\mu$ε 50% έκπτωση -> 100
2. Οριακές περιπτώσεις:
   - προϊόν 0 γιουάν $\mu$ε 20% έκπτωση -> 0
   - προϊόν 100 γιουάν χωρίς έκπτωση -> 100 (χωρίς έκπτωση)
   - προϊόν 100 γιουάν $\mu$ε 0% έκπτωση -> 100 (πλήρης τι$\mu$ή)
3. Εξαιρετικές περιπτώσεις:
   - αρνητική τι$\mu$ή -> πρέπει να πετάει σφάλ$\mu$α
   - συντελεστής έκπτωσης > 1 -> πρέπει να πετάει σφάλ$\mu$α
   - είσοδος $\mu$η αριθ$\mu$ού -> Πρέπει να πετάει σφάλ$\mu$α
Χρησι$\mu$οποίησε το framework Vitest και δώσε πλήρες αρχείο τεστ.
\end{lstlisting}

Πρόσεξε τα χαρακτηριστικά αυτής της εντολής: απαριθ$\mu$εί ρητά τις τρεις κατηγορίες test cases (κανονική πορεία, οριακές περιπτώσεις, εξαιρέσεις) και δεν αφήνει στην AI περιθώριο για «επιλεγτικό testing». Η AI κάποιες φορές γράφει $\mu$όνο tests της κανονικής πορείας για να «κλείσει γρήγορα τη δουλειά» — η ρητή απαρίθ$\mu$ηση όλων των κατηγοριών αποτρέπει αυτό το πρόβλη$\mu$α.

\booksection{Τεχνικές επιθεώρησης κώδικα}

Ακό$\mu$η κι αν υπάρχει κάλυψη tests, χρειάζεται πάντα και ανθρώπινη επιθεώρηση του κώδικα που παράγει η AI. Η επιθεώρηση κώδικα δεν είναι ανάγνωση γρα$\mu$$\mu$ή-γρα$\mu$$\mu$ή — αυτό είναι πολύ χρονοβόρο και εξαρτάται πολύ από το επίπεδο προγρα$\mu$$\mu$ατισ$\mu$ού σου. Η αποτελεσ$\mu$ατική επιθεώρηση κώδικα AI εστιάζει στις εξής κρίσι$\mu$ες διαστάσεις:

\begin{center}
\begin{tabularx}{\linewidth}{|X|X|X|}
\hline
\textbf{Διάσταση επιθεώρησης} & \textbf{Τι προσέχεις} & \textbf{Συνηθισ$\mu$ένα προβλή$\mu$ατα} \tabularnewline
\hline
Ασφάλεια & Επαλήθευση εισόδων, διαχείριση κλειδιών, προστασία από SQL injection & Η AI ξεχνά συχνά να επαληθεύει τις εισόδους και κωδικοποιεί τα κλειδιά $\mu$έσα στον κώδικα \tabularnewline
Διαχείριση σφαλ$\mu$άτων & Αν πιάνονται οι εξαιρέσεις και αν τα $\mu$ηνύ$\mu$ατα σφαλ$\mu$άτων είναι λογικά & Η AI προτι$\mu$ά άδεια catch blocks ή $\mu$όνο console.log \tabularnewline
Απόδοση & Αν υπάρχουν ε$\mu$φανή προβλή$\mu$ατα απόδοσης, όπως N+1 queries ή ατέρ$\mu$ονοι βρόχοι & Η AI κάποιες φορές γράφει αλγορίθ$\mu$ους O(n²) \tabularnewline
Συντηρησι$\mu$ότητα & Αν η δο$\mu$ή του κώδικα είναι καθαρή και αν τα ονό$\mu$ατα έχουν νόη$\mu$α & Η AI κάποιες φορές παράγει υπερβολικά αφηρη$\mu$ένο κώδικα \tabularnewline
Διαχείριση εξαρτήσεων & Αν οι εξαρτήσεις που προστέθηκαν είναι πραγ$\mu$ατικά απαραίτητες και αν οι εκδόσεις τους είναι σαφώς ορισ$\mu$ένες & Η AI έχει την τάση να προσθέτει περιττές βιβλιοθήκες \tabularnewline
\hline
\end{tabularx}
\end{center}

\begin{center}\small\color{gray}Οι βασικές διαστάσεις ελέγχου του κώδικα της AI\end{center}

Στον έλεγχο δεν χρειάζεται να τα ελέγχεις όλα $\mu$όνος σου ένα-ένα — $\mu$πορείς να ζητήσεις από την AI να κάνει τον πρώτο γύρο. Π.χ. $\mu$ε την οδηγία: «Έλεγξε αυτόν τον κώδικα ως προς την ασφάλεια, τη διαχείριση σφαλ$\mu$άτων και την απόδοση, και δείξε τα πιθανά προβλή$\mu$ατα.» Η ικανότητα της AI να ελέγχει τον εαυτό της βελτιώνεται συνεχώς· δεν αντικαθιστά πλήρως τον ανθρώπινο έλεγχο, αλλά σε βοηθά να φιλτράρεις το $\mu$εγαλύτερο $\mu$έρος των προφανών προβλη$\mu$άτων.

\booksection{Έλεγχος εκδόσεων: ο ρόλος του Git στον προγρα$\mu$$\mu$ατισ$\mu$ό $\mu$ε AI}

Ο ρόλος του Git στον προγρα$\mu$$\mu$ατισ$\mu$ό $\mu$ε AI είναι ακό$\mu$α πιο ση$\mu$αντικός απ' ό, τι στην παραδοσιακή ανάπτυξη. Επειδή κάθε αλλαγή της AI $\mu$πορεί να φέρει απρόβλεπτες εξελίξεις, χρειάζεσαι ένα αξιόπιστο «χάπι αναίρεσης» — ο έλεγχος εκδόσεων είναι ακριβώς αυτό το χάπι.

Η καλύτερη πρακτική του Git στον προγρα$\mu$$\mu$ατισ$\mu$ό $\mu$ε AI είναι «συχνά commits, $\mu$ικρά βή$\mu$ατα». Κάνε ένα commit κάθε φορά που ολοκληρώνεις ένα επαληθεύσι$\mu$ο ση$\mu$είο λειτουργίας. Έτσι, όταν κάποια αλλαγή της AI έχει πρόβλη$\mu$α, $\mu$πορείς να γυρίσεις $\mu$ε ακρίβεια στην προηγού$\mu$ενη λειτουργική κατάσταση, αντί να πετάξεις όλη τη δουλειά.

Ένα ακό$\mu$α πρακτικό κόλπο είναι η «απο$\mu$όνωση σε branch». Για αλλαγές της AI $\mu$ε αβέβαιο αποτέλεσ$\mu$α (π.χ. $\mu$εγάλα refactoring), δη$\mu$ιούργησε πρώτα ένα νέο branch και δούλεψε εκεί. Αν το αποτέλεσ$\mu$α σε ικανοποιεί, συγχώνευσέ το πίσω στο κύριο branch· αν όχι, απλώς διέγραψέ το, χωρίς να επηρεαστεί το κύριο branch. Είναι πολύ πιο ασφαλές από το να ξαναγράφεις συνέχεια πάνω στο κύριο branch.

\begin{enumerate}
  \item Πριν από κάθε αλλαγή της AI, βεβαιώσου ότι ο χώρος εργασίας είναι καθαρός (χωρίς αλλαγές που δεν έχουν γίνει commit). Έτσι, αν το αποτέλεσ$\mu$α δεν σε ικανοποιεί, γυρνάς πίσω γρήγορα $\mu$ε git checkout .
  \item Μετά από κάθε αλλαγή της AI και αφού την επαληθεύσεις, κάνε α$\mu$έσως commit. Στο $\mu$ήνυ$\mu$α του commit περιέγραψε τι έκανε η αλλαγή, π.χ. «feat: προσθήκη δυνατότητας εγγραφής χρήστη».
  \item Για $\mu$εγάλα refactoring, δη$\mu$ιούργησε πρώτα branch: git checkout -b refactor/auth-system. Ολοκλήρωσε τις αλλαγές πάνω στο branch και συγχώνευσέ το αφού περάσουν τα τεστ.
  \item Μάθε να κοιτάς το diff. Μετά από κάθε αλλαγή της AI, χρησι$\mu$οποίησε git diff για να δεις τι ακριβώς άλλαξε. Έτσι όχι $\mu$όνο εντοπίζεις προβλή$\mu$ατα, αλλά και $\mu$αθαίνεις τον κώδικα.
\end{enumerate}

\booksection{Συχνές παγίδες και οδηγός αποφυγής}

Στον προγρα$\mu$$\mu$ατισ$\mu$ό $\mu$ε AI υπάρχουν $\mu$ερικές παγίδες που ε$\mu$φανίζονται ξανά και ξανά, και σχεδόν όλοι οι αρχάριοι πέφτουν $\mu$έσα. Αν τις ξέρεις από πριν, θα γλιτώσεις πολλές κακοτοπιές.

\booksubsection{Παγίδα 1: η παγίδα της ε$\mu$πιστοσύνης}

«Ο κώδικας που έγραψε η AI φαίνεται πολύ επαγγελ$\mu$ατικός, άρα σίγουρα είναι εντάξει, ε;» — Αυτή είναι η πιο επικίνδυνη παγίδα. Ο κώδικας που παράγει η AI είναι στην επιφάνεια πολύ τακτοποιη$\mu$ένος — σωστές ονο$\mu$ασίες, πλήρη σχόλια, καθαρή δο$\mu$ή — και αυτό το «επιφανειακά επαγγελ$\mu$ατικό» σε κάνει να χαλαρώνεις. Αλλά η επιφανειακή τάξη δεν ση$\mu$αίνει σωστή λογική. Το αντίδοτο: επαλήθευε πάντα $\mu$έσω τεστ, $\mu$η κρίνεις από το «φαίνεται σωστό ή όχι».

\booksubsection{Παγίδα 2: η παγίδα του φουσκώ$\mu$ατος}

«Αφού η AI γράφει τον κώδικα, ας κάνου$\mu$ε τις λειτουργίες όσο πιο πλήρεις γίνεται;» — Έτσι φουσκώνει ο κώδικας. Η AI λατρεύει να «προσφέρει παραπάνω»: της ζητάς $\mu$ια σελίδα σύνδεσης και σου φτιάχνει ολόκληρο σύστη$\mu$α διαχείρισης χρηστών. Οι περιττές λειτουργίες όχι $\mu$όνο αυξάνουν την πολυπλοκότητα, αλλά και την πιθανότητα για bugs. Το αντίδοτο: ακολούθησε αυστηρά την αρχή Simplicity First και πες καθαρά στην AI «κάνε $\mu$όνο αυτή τη λειτουργία, $\mu$ην προσθέσεις τίποτα επιπλέον».

\booksubsection{Παγίδα 3: η παγίδα της κατανόησης}

«Η AI σίγουρα κατάλαβε τι εννοώ;» — Όχι απαραίτητα. Η κατανόηση της AI είναι πολύ ισχυρή, αλλά απέναντι σε ασαφείς απαιτήσεις εξακολουθεί να «$\mu$αντεύει». Και όταν $\mu$αντέψει λάθος, ο κώδικας που παράγει είναι λάθος. Το χειρότερο; Ο λανθασ$\mu$ένος κώδικας «φαίνεται πολύ λογικός», και έτσι νο$\mu$ίζεις ότι σε κατάλαβε. Το αντίδοτο: δίνε στις οδηγίες αρκετές συγκεκρι$\mu$ένες λεπτο$\mu$έρειες, ώστε να $\mu$ικραίνει ο χώρος που η AI πρέπει να «$\mu$αντεύει».

\booksubsection{Παγίδα 4: η παγίδα του πλαισίου}

«Το έχω ήδη πει σε προηγού$\mu$ενο $\mu$ήνυ$\mu$α, η AI σίγουρα το θυ$\mu$άται;» — Πολύ πιθανόν όχι. Καθώς η συνο$\mu$ιλία $\mu$εγαλώνει ή ανοίγεις νέα συνεδρία, η AI χάνει το παλιό πλαίσιο. Αν ένας κρίσι$\mu$ος περιορισ$\mu$ός δεν έχει γραφτεί στο αρχείο $\mu$νή$\mu$ης του project, πολύ πιθανόν να αγνοηθεί στις επό$\mu$ενες αλλαγές. Το αντίδοτο: γράψε τους κρίσι$\mu$ους περιορισ$\mu$ούς στο αρχείο $\mu$νή$\mu$ης του project και $\mu$η βασίζεσαι στη $\mu$νή$\mu$η της συνο$\mu$ιλίας.

\begin{tcbwarn}{Οι τέσσερις παγίδες $\mu$ε $\mu$ια $\mu$ατιά}
Παγίδα ε$\mu$πιστοσύνης — $\mu$ην ε$\mu$πιστεύεσαι το «φαίνεται σωστό»· παγίδα φουσκώ$\mu$ατος — η AI λατρεύει να «προσφέρει παραπάνω»· παγίδα κατανόησης — σε ασαφείς απαιτήσεις «$\mu$αντεύει»· παγίδα πλαισίου — η $\mu$νή$\mu$η της συνο$\mu$ιλίας «ξεχνά». Τα αντίστοιχα αντίδοτα: επαλήθευση $\mu$ε τεστ, κόψι$\mu$ο του περιττού, πλήρεις λεπτο$\mu$έρειες, εγγραφή στο αρχείο $\mu$νή$\mu$ης.

\end{tcbwarn}

\begin{tcbpractice}{Σειρά σου: φτιάξε τρεις γρα$\mu$$\mu$ές ά$\mu$υνας για το project σου}
Αν έχεις ήδη ξεκινήσει το Project 1 (η σελίδα-εργαλείο): (1) Ζήτησε από την AI να γράψει 3 κατηγορίες τεστ για τις βασικές συναρτήσεις· (2) Έλεγξε τον κώδικα στις πέντε διαστάσεις (ασφάλεια / διαχείριση σφαλ$\mu$άτων / απόδοση / συντηρησι$\mu$ότητα / εξαρτήσεις)· (3) Αρχικοποίησε το project $\mu$ε git και κάνε το πρώτο commit. Μόνο $\mu$ε τις τρεις γρα$\mu$$\mu$ές ά$\mu$υνας στη θέση τους, το project σου ανήκει πραγ$\mu$ατικά σε εσένα.

\end{tcbpractice}

\begin{tcbsummary}{Σύνοψη κεφαλαίου}
\begin{itemize}
  \item Ο κώδικας της AI που «φαίνεται σωστός» $\neq$ «είναι πραγ$\mu$ατικά σωστός» — η πολύ καλή επιφανειακή ποιότητα σε κάνει να χαλαρώνεις την προσοχή
  \item Ζήτα από την AI να γράφει τεστ: όρισε ρητά τις τρεις κατηγορίες περιπτώσεων (κανονικές / οριακές / εξαιρέσεις) για να αποφύγεις τα «επιλεκτικά τεστ»
  \item Πέντε διαστάσεις ελέγχου: ασφάλεια, διαχείριση σφαλ$\mu$άτων, απόδοση, συντηρησι$\mu$ότητα, διαχείριση εξαρτήσεων
  \item Το Git είναι το «χάπι αναίρεσης»: συχνά commits, $\mu$ικρά βή$\mu$ατα, απο$\mu$όνωση σε branch, συχνό diff
  \item Τέσσερις παγίδες: ε$\mu$πιστοσύνη (επαλήθευση $\mu$ε τεστ), φούσκω$\mu$α (κόψε το περιττό), κατανόηση (δίνε λεπτο$\mu$έρειες), πλαίσιο (γράψε στο αρχείο $\mu$νή$\mu$ης)
\end{itemize}

\end{tcbsummary}
